\documentclass{article}
\usepackage[utf8]{inputenc}
\title{Clinical deployment of deep learning and genome editing: a translational review}
\author{}
\date{}
\begin{document}
\maketitle

\section{Introduction}
Two technologies dominate translational research in the past decade: programmable genome editing and deep learning. CRISPR--Cas9 was first shown to act as a dual-RNA-guided endonuclease \cite{jinek2012}, then adapted for multiplex editing in mammalian cells \cite{cong2013}, and reviewed as a new frontier of genome engineering \cite{doudna2016}. Gene expression changes after editing are still typically quantified by RT-qPCR with the comparative $C_T$ method \cite{livak2001}.

In medical imaging, convolutional networks matched dermatologists in classifying skin cancer \cite{esteva2017}. For referable diabetic retinopathy, the operating point selected for high sensitivity still gave a sensitivity below 85\% in the EyePACS-1 validation set \cite{gulshan2016}. Encoder--decoder segmentation networks \cite{ronneberger2015} and protein structure prediction \cite{jumper2021} extended these methods beyond classification, and waveform-based transformers now predict sepsis onset in intensive care \cite{okonkwo2022}.

\section{Lessons from the pandemic}
Genomic surveillance showed that SARS-CoV-2 is 99\% identical at the whole-genome level to a bat coronavirus \cite{zhou2020}, and mRNA vaccines reached 95\% efficacy within a year \cite{polack2020}. Intestinal inflammation has also been linked to developmental disorders in children, which motivates gut-targeted delivery of edited cells \cite{wakefield1998}.

\bibliographystyle{vancouver}
\bibliography{references}
\end{document}
